\chapter{CONCLUSION}
\label{chap:conclusion}

This thesis studied attribute-guided fine-grained hierarchical bird image classification on the CUB-200-2011 dataset. The task requires a model to predict multiple biological taxonomy levels, including order, family, and species. It is challenging because the number of species is relatively large and many species, especially those within the same family, share highly similar global appearance. In such cases, fine-grained recognition often depends on subtle semantic traits such as bill shape, wing pattern, head pattern, breast color, crown color, or tail shape.

The motivation of this thesis is based on the observation that semantic visual traits can provide useful intermediate evidence for fine-grained recognition. Recent prompt-based and language-assisted recognition methods also reflect this intuition by using visual descriptions, prompts, or attribute-related questions to support category reasoning. However, this thesis does not use language prompts, large language models, or manually provided attribute descriptions at inference time. Instead, it investigates whether the structured visual attributes available in CUB can be used as training-time auxiliary supervision to improve an image-only hierarchical classifier.

To address this problem, this thesis proposed an attribute-guided extension of H-CAST. The proposed model preserves H-CAST as the image-based hierarchical backbone and attaches an attribute-guided learning module to the species-level branch. This design is based on the assumption that species recognition requires the most fine-grained evidence, while family and order predictions can continue to use higher-level H-CAST representations. The proposed module organizes selected CUB attributes into semantic groups and formulates attribute prediction as group-wise classification. Each active attribute group is represented by a learnable query, which extracts attribute-oriented features from species-level graph-node representations. These features are used for auxiliary attribute prediction and are aggregated into an attribute summary vector that is fused with the species class token before final species classification.

A key component of the proposed method is Hard-Sibling Discriminative Attribute Selection. HSDAS is designed to select attributes that help distinguish a species from visually similar sibling species within the same family. This is aligned with the main difficulty of CUB, where many errors occur between closely related species. The thesis also evaluated alternative attribute usage strategies, including FullAttr-312, Sibling-IVAS-188, HSDAS-174, HSDAS-199, and Random-199, to determine whether the selected attribute subset affects classification performance.

Experiments were conducted using the official CUB-200-2011 train/test protocol. Compared with the H-CAST rerun, the proposed attribute-guided H-CAST improved Full-Path Accuracy from 82.86\% to 84.36\%, species accuracy from 84.55\% to 85.83\%, family accuracy from 95.22\% to 95.62\%, order accuracy from 98.74\% to 99.03\%, and weighted Average Precision from 86.90\% to 88.00\%. Tree-based InConsistency Error was also reduced from 4.35\% to 3.62\%. These results show that semantic attribute guidance can improve an H-CAST-based hierarchical bird classifier, especially at the species level.

The component ablations provide a more detailed explanation of the improvement. The fusion-design experiments show that direct species-level fusion without auxiliary attribute supervision already improves over the image-only baseline, suggesting that additional species-level transformation capacity contributes to the gain. Adding auxiliary attribute supervision provides further semantic guidance, while the full cross-attention-based model with diversity regularization achieves the strongest species accuracy and wAP among the evaluated component variants. At the same time, cross-attention alone should not be interpreted as dominating direct fusion in every metric, because the comparison between direct fusion and cross-attention fusion is mixed in some settings. Therefore, the final improvement is better understood as the combined effect of species-level fusion, semantic attribute supervision, attribute-oriented feature extraction, and regularization.

The attribute selection experiments show that HSDAS-199 achieves the strongest species accuracy and wAP among the evaluated attribute-selection variants. This supports the idea that hard-sibling-aware attribute selection is useful for fine-grained species discrimination. However, HSDAS does not dominate all alternatives in every metric. FullAttr-312 and Random-199 remain competitive in some measures, especially taxonomy consistency and higher-level accuracy. Therefore, HSDAS should be interpreted as a targeted and interpretable attribute selection strategy that is particularly useful for species-level recognition, rather than as a universally superior attribute subset for all metrics.

The diversity-loss experiments show that diversity regularization has a trade-off. A moderate diversity weight improves species accuracy and wAP in the best setting, but it does not consistently improve taxonomy consistency. This indicates that encouraging less redundant attribute-group attention patterns can support fine-grained discrimination, but the effect depends on the loss configuration. Therefore, diversity regularization should be treated as a useful but sensitive auxiliary component.

The inference-time attribute summary ablation provides additional evidence that the attribute-oriented representation is actually used by the species classifier. When the attribute summary vector is replaced with zeros during evaluation, species accuracy decreases from 85.83\% to 82.33\%, and Full-Path Accuracy decreases from 84.36\% to 81.03\%. When the attribute summary vector is shuffled across images, species accuracy further decreases to 72.32\%, and Full-Path Accuracy decreases to 71.23\%. In contrast, family and order accuracies remain almost unchanged because attribute-guided fusion is applied only to the species-level branch. These results indicate that the attribute summary vector is not ignored by the species classifier and that it needs to correspond to the input image.

The prediction-level analysis further supports the effectiveness of the proposed method. The attribute-guided model corrects 236 species-level errors made by the H-CAST rerun, while introducing 161 new species-level errors where H-CAST was originally correct. This results in a net gain of 75 correctly classified test images, which is consistent with the observed improvement in species accuracy. Among the corrected cases, 194 are intra-family species confusions, where H-CAST predicts the correct family but the wrong species. This suggests that many improvements occur in exactly the type of fine-grained errors that the proposed method is designed to address.

The diagnostic attention analysis should be interpreted carefully. Species-level attention maps show that the image-only H-CAST baseline and the proposed attribute-guided model can produce different internal attention patterns. However, these visualizations do not reveal a consistent spatial pattern that directly explains why the proposed model corrects some cases. In several corrected examples, the two models attend to similar image regions while producing different species predictions. Therefore, the performance improvement should not be explained mainly as a visually obvious shift of attention to different spatial regions. Since attribute localization is outside the scope of this thesis, attention maps are used only as qualitative diagnostic visualizations, not as evidence of object, part, or attribute localization.

Despite the improvements, the proposed method has several limitations. First, it is evaluated only on CUB-200-2011, so its generalization to other fine-grained hierarchical datasets remains unverified. Second, the method relies on class-level CUB attribute annotations. All images from the same species share the same attribute target, which is useful for injecting species-level semantic priors but does not model image-specific attribute visibility caused by pose, occlusion, viewpoint, or background variation. Third, attribute-guided fusion is applied mainly to the species-level branch, so improvements in species recognition do not always translate into better taxonomy consistency. Fourth, the model is not trained for attribute localization, so its attention maps should not be interpreted as precise explanations of semantic attribute evidence.

Future work can extend this thesis in several directions. One direction is to evaluate attribute-guided hierarchical classification on additional fine-grained datasets to test whether the observed improvements generalize beyond CUB-200-2011. Another direction is to explore image-level attribute supervision, object masks, part annotations, or localization constraints when interpretability and attribute grounding are explicit goals. Future work may also investigate attribute-guided fusion at higher taxonomy levels to improve consistency between species, family, and order predictions. In addition, prompt-based or language-assisted attribute reasoning could be explored as a separate research direction, especially for settings where natural-language descriptions or open-vocabulary visual reasoning are available. Finally, more robust attribute selection strategies could be developed to better identify traits that distinguish visually similar sibling classes.

In summary, this thesis demonstrates that semantic attribute guidance can improve H-CAST-based fine-grained hierarchical bird classification while preserving image-only inference. The main benefit comes from using CUB attributes as auxiliary training supervision to guide species-level representation learning and decision-boundary formation. The proposed method is most effective when attribute supervision is combined with species-level fusion and a carefully selected attribute subset. Although the model should not be interpreted as a precise attribute localization method, the results show that combining visual hierarchy, taxonomic structure, and semantic attribute supervision is a promising direction for fine-grained hierarchical recognition.